% =========================================================
%  sections/09-evaluasi-refleksi.tex
%  Bagian 9 — Evaluasi dan Refleksi Pembelajaran
% =========================================================

\section{Evaluasi dan Refleksi Pembelajaran}

\subsection{Penilaian Formatif}

Penilaian formatif dilaksanakan sepanjang semester untuk memantau kemajuan
mahasiswa dan memberikan umpan balik sebelum asesmen sumatif.

\begin{itemize}
  \item \textbf{Umpan balik praktikum mingguan.} Dosen dan asisten memeriksa
        artefak praktikum setiap minggu dan memberikan catatan perbaikan
        langsung pada repositori mahasiswa.
  \item \textbf{Kuis singkat.} Dua kuis pada Minggu 5 dan Minggu 8 digunakan
        untuk mendeteksi miskonsepsi lebih awal, khususnya pada konsep
        \istilah{delta time}, resolusi \istilah{collision}, dan transisi FSM.
  \item \textbf{Umpan balik sejawat.} Pada Minggu 9 dan Minggu 15 mahasiswa
        saling menilai hasil kerja kelompok lain menggunakan lembar
        \istilah{playtest} terstruktur, kemudian mendiskusikan temuannya.
  \item \textbf{Demonstrasi kelas.} Setiap kelompok memperagakan kemajuan
        proyeknya secara berkala agar memperoleh masukan sebelum tenggat akhir.
\end{itemize}

\subsection{Penilaian Sumatif}

Penilaian sumatif diarahkan untuk menilai ketercapaian CPMK secara menyeluruh
melalui UTS (Minggu 10) dan UAS (Minggu 16), serta melalui penilaian mini
proyek AI (Minggu 9) dan proyek akhir (Minggu 15--16). Bobot dan
keterkaitannya dengan CPMK diuraikan pada Bagian~\ref{sec:penilaian}.

\subsection{Refleksi Mahasiswa}

Setiap mahasiswa mengisi lembar refleksi pada akhir semester yang memuat tiga
pertanyaan pokok:

\begin{enumerate}
  \item Konsep apa yang paling sulit dipahami pada blok Python/Pygame dan blok
        Unity, dan bagaimana usaha yang dilakukan untuk mengatasinya?
  \item Apa perbedaan pengalaman bekerja dengan Python/Pygame dibandingkan
        dengan Unity, ditinjau dari kemudahan, kendali atas algoritma, dan
        kecepatan pengembangan?
  \item Kompetensi apa yang ingin dikembangkan setelah mata kuliah ini, dan
        bagaimana rencana untuk mencapainya?
\end{enumerate}

Hasil refleksi dibahas dalam sesi penutup pada Minggu 15 dan direkapitulasi
oleh dosen sebagai masukan perbaikan pelaksanaan perkuliahan.

\subsection{Refleksi Dosen dan Penjaminan Mutu}

Setelah perkuliahan berakhir, dosen melakukan analisis ketercapaian CPMK
dengan menghitung persentase mahasiswa yang mencapai ambang batas pada setiap
CPMK. Analisis tersebut dilaporkan kepada program studi sebagai bagian dari
siklus penjaminan mutu (\istilah{Plan--Do--Check--Act}) dan menjadi dasar
penyesuaian RPS pada semester berikutnya.

\begin{table}[htbp]
\centering
\caption{Ambang ketercapaian CPMK dan tindak lanjut}
\label{tab:ambang}
\begin{tabularx}{\textwidth}{@{}C{1.8cm} L{3.2cm} C{2.2cm} L{3.2cm} Y@{}}
\toprule
\headrow \textbf{CPMK} & \textbf{Indikator Ketercapaian} & \textbf{Ambang Batas} & \textbf{Instrumen Pemantauan} & \textbf{Tindak Lanjut Bila Belum Tercapai} \\
\midrule
\CPMK{1} &
Mahasiswa mampu merancang dan mengimplementasikan \istilah{game loop}, gerak berbasis \istilah{delta time}, fisika sederhana, dan \istilah{collision}. &
$\geq 75\%$ mahasiswa bernilai $\geq 70$ &
Praktikum 01--05, Kuis 1, UTS &
Menyediakan sesi remedial praktikum dan contoh kode pembanding; memperjelas indikator pada RPS berikutnya. \\
\addlinespace
\CPMK{2} &
Mahasiswa mampu mengimplementasikan perilaku \istilah{chase}, FSM, dan A* serta memadukannya. &
$\geq 75\%$ mahasiswa bernilai $\geq 70$ &
Praktikum 06--09, Kuis 2, Mini Proyek AI &
Menambah latihan terarah pada bagian yang paling lemah (FSM atau A*) dan memperpanjang bimbingan proyek. \\
\addlinespace
\CPMK{3} &
Mahasiswa mampu membangun level Unity yang dapat dimainkan dengan fisika yang benar. &
$\geq 75\%$ mahasiswa bernilai $\geq 70$ &
Praktikum 09--10, Proyek Akhir &
Menambah demonstrasi terbimbing dan menyediakan proyek contoh sebagai acuan tahapan. \\
\addlinespace
\CPMK{4} &
Mahasiswa mampu mengintegrasikan aset visual, UI, dan audio serta mengevaluasi prototipe. &
$\geq 75\%$ mahasiswa bernilai $\geq 70$ &
Praktikum 11--12, Proyek Akhir, UAS &
Menyediakan pustaka aset siap pakai dan sesi \istilah{playtest} tambahan sebelum tenggat akhir. \\
\bottomrule
\end{tabularx}
\end{table}

Apabila suatu CPMK belum mencapai ambang batas, dosen mencatat penyebabnya
(materi, metode, atau asesmen) dan mengusulkan perbaikan yang akan diterapkan
pada pelaksanaan semester berikutnya.
